\documentclass[a4paper,11pt]{article}
\usepackage[utf8]{inputenc}
\usepackage[T1]{fontenc}
\usepackage[french]{babel}
\usepackage{lmodern}
\usepackage{amsmath,amssymb,amsthm}
\usepackage{mathrsfs}
\usepackage{enumitem}
\usepackage{multicol}

\setlist[enumerate]{label=\textbf{\textcolor{Couleur8}{\arabic*.}}, left=1em}

% Si vous voulez aussi modifier les listes enumerate imbriquées :
\setlist[enumerate,2]{label=\textcolor{Couleur8}{\alph*)}}
\setlist[enumerate,3]{label=\textcolor{Couleur8}{\roman*)}}
\setlist[enumerate,4]{label=\textcolor{Couleur8}{\Alph*)}}
\usepackage{graphicx}
\usepackage{xcolor}
\usepackage{tcolorbox}
\usepackage{geometry}
\usepackage{fancyhdr}
\usepackage{titlesec}
\usepackage{cancel}
\usepackage{ifthen}
\usepackage{tikz}
\usetikzlibrary{arrows.meta}
\usepackage{tkz-tab}
\usepackage{eurosym}

% OPTION PROF/ÉLÈVE
% Décommentez UNE des deux lignes suivantes :
\newboolean{prof}\setboolean{prof}{true}   % Version professeur (avec corrections des devoirs)
\newboolean{prof}\setboolean{prof}{true}  % Version élève (sans corrections des devoirs)

\definecolor{Couleur1}{HTML}{4A5D7A} %Th
\definecolor{Couleur2}{HTML}{A5B5D6} %Prop
\definecolor{Couleur3}{HTML}{A8C68A} %Méthode
\definecolor{Couleur6}{HTML}{8A9CC1} %Def
\definecolor{Couleur7}{HTML}{E6B459} %Rq Voc
\definecolor{Couleur8}{HTML}{D36740} %Titres
\definecolor{correction}{HTML}{D36740} % Couleur pour les corrections
\definecolor{coopmathscorr}{HTML}{F15929} % Couleur corrections coopmaths

% Configuration des marges
\geometry{
	top=2cm,
	bottom=2cm,
	inner=1.5cm,
	outer=1.5cm,
}

% Police
\renewcommand{\familydefault}{\sfdefault}

% Compteurs
\newcounter{seancenum}
\newcounter{seancenumD}
\setcounter{seancenum}{1}
\setcounter{seancenumD}{1}

% Configuration des en-têtes et pieds de page
\pagestyle{fancy}
\fancyhf{}
\ifthenelse{\boolean{prof}}{
    \fancyhead[L]{\textcolor{Couleur8}{\textbf{Corrections Automatismes - Classe de seconde - Version Professeur}}}
}{
    \fancyhead[L]{\textcolor{Couleur8}{\textbf{Corrections Devoirs - Séance 3 - Classe de seconde}}}
}
\fancyhead[R]{\textcolor{Couleur8}{\textbf{2025-2026}}}
\fancyfoot[L]{\textcolor{Couleur8}{Mme Bertail et Mme Pinto}}
\fancyfoot[R]{\textcolor{Couleur8}{\thepage}}
\renewcommand{\headrulewidth}{1pt}
\renewcommand{\headrule}{\hbox to\headwidth{\color{Couleur8}\leaders\hrule height \headrulewidth\hfill}}

% Environnements personnalisés
\newtcolorbox{seance}[1]{
    colback=Couleur6!10,
    colframe=Couleur1!65,
    fonttitle=\bfseries\large,
    title=Correction Séance \theseancenum #1,
    left=5mm,
    right=5mm,
    top=3mm,
    bottom=3mm,
    arc=0mm,
    after=\stepcounter{seancenum}
}

\newtcolorbox{seanceD}[1]{
    colback=Couleur6!5,
    colframe=Couleur6!45,
    fonttitle=\bfseries\large,
    title=Correction Devoirs - Séance \theseancenumD #1,
    left=5mm,
    right=5mm,
    top=3mm,
    bottom=3mm,
    arc=0mm,
    after=\stepcounter{seancenumD}
}

\begin{document}

\setcounter{seancenumD}{3}
\newpage
\begin{seanceD}{} %3

\begin{enumerate}
\item De la relation $-4x-5y=8$, on déduit en ajoutant $4x$ dans chaque membre :
$-5y=8+4x$
Puis, en divisant par $-5$, on obtient : $y=\dfrac{8+4x}{-5}$ ou ${\color{correction}\boldsymbol{y=\dfrac{-8-4x}{5}}}$.


\item  On met l'expression au même dénominateur : \\$
        \dfrac{1}{7}-\dfrac{7x+5}{x}=\dfrac{x-7\times \left(7x+5\right)}{7x}=\dfrac{x -49x -35}{7x}=\dfrac{-48x -35}{7x}=-\dfrac{48x +35}{7x}$\\La bonne réponse est la réponse {\bfseries \color{correction}D}.

\item On remplace $a$, $b$, $c$ et $d$ par les valeurs données : \\
$F=\dfrac{1}{7}+\dfrac{2}{-5\times \dfrac{1}{5}}=\dfrac{1}{7}+\dfrac{2}{-1}=\dfrac{1}{7}-2={\color{correction}\boldsymbol{-\dfrac{13}{7}}}$

\item Un terme de la suite peut s'exprimer sous la forme $2\times n-1$ avec $n$ le rang du terme.
Par exemple, le $4^{\text{e}}$ terme est $2\times 4-1=7$.
Le $17^{\text{e}}$ nombre est donc $2\times 17-1={\color{correction}\boldsymbol{33}}$.

\item Il y a ${\color{correction}\boldsymbol{6}}$ entiers dans l'intervalle $\bigg[-3{,}4\, ; \,2\bigg]$.
Mentalement : Comptez-les ! Ce sont : $-3, -2, -1, 0, 1, 2$.

\item L'intersection $]-\infty\,;\, 8[\,\cap \,]-\infty\,;\,9]$ = ${\color{correction}\boldsymbol{]-\infty\,;\,8[}}$.
Les nombres de l'intervalle $]-\infty\,;\,8[$ appartiennent à l'intervalle $]-\infty\,;\, 8[$ et à l'intervalle $]-\infty\,;\,9]$.
\end{enumerate}

\end{seanceD}

\end{document}